% =====================================================================
\chapter{Modulo 4 -- Classi e Oggetti: Costruttori, Distruttore, Regola dei Tre}
% =====================================================================

% ---------------------------------------------------------------------
\section{Teoria \& Analogia Chiave}
% ---------------------------------------------------------------------

Nel Modulo 1 abbiamo visto \emph{cosa} sono classi e oggetti. In questo modulo
vediamo \emph{come vivono}: come nascono, come muoiono, come si copiano. È il
modulo che collega la OOP alla memoria dinamica del Modulo 3. Quando un oggetto
possiede memoria nello heap, nascita, morte e copia vanno gestite a mano, ed è
qui che si concentrano gli errori più gravi (e le domande d'orale più frequenti).

\subsection*{Il ciclo di vita di un oggetto}

Ogni oggetto attraversa tre fasi:
\begin{enumerate}
  \item \textbf{nasce}: viene riservata la memoria e viene eseguito
        automaticamente il \textbf{costruttore};
  \item \textbf{vive}: riceve messaggi, cambia stato;
  \item \textbf{muore}: viene eseguito automaticamente il \textbf{distruttore},
        poi la memoria viene restituita.
\end{enumerate}
\emph{Quando} avvengono nascita e morte dipende da dove l'oggetto è stato creato,
con le stesse regole delle variabili del Modulo 3:
\begin{center}\small
\begin{tabular}{p{3.4cm}p{5.4cm}p{5.4cm}}
\toprule
\textbf{Oggetto} & \textbf{Nasce} & \textbf{Muore} \\
\midrule
globale & prima del \cc|main| & dopo la fine del \cc|main| \\
locale (stack) & quando l'esecuzione arriva alla dichiarazione & alla \cc|}| del suo blocco \\
dinamico (heap) & con \cc|new| & con \cc|delete| (se te ne dimentichi: mai) \\
parametro per valore & all'ingresso nella funzione (è una copia) & all'uscita dalla funzione \\
\bottomrule
\end{tabular}
\end{center}
Gli oggetti locali di uno stesso blocco muoiono in \textbf{ordine inverso} rispetto
a come sono nati: l'ultimo creato è il primo distrutto, come i piatti di una pila.
L'Esercizio~4.2 lo mostra con dei messaggi stampati da costruttore e distruttore.

\subsection*{Il costruttore}

\begin{definizione}[: costruttore]
Il \textbf{costruttore} è un metodo speciale che viene eseguito
\textbf{automaticamente} quando un oggetto nasce. Ha \textbf{lo stesso nome della
classe}, \textbf{nessun tipo di ritorno} (nemmeno \cc|void|) e può avere
parametri. Il suo compito è portare l'oggetto in uno stato iniziale valido.
\end{definizione}

\begin{analogia}[: il collaudo in fabbrica]
Un'auto non esce dalla fabbrica con il motore smontato e il serbatoio pieno di
sabbia: prima della consegna c'è il collaudo, che la mette in condizioni di
funzionare. Il costruttore è il collaudo: quando il programma riceve l'oggetto,
questo è già pronto e coerente.
\end{analogia}

Una classe può avere \textbf{più costruttori} (overloading), purché abbiano
parametri diversi per numero o tipo. Il compilatore sceglie quello giusto in base
agli argomenti che scrivi:
\begin{lstlisting}
class Rettangolo {
    double base, altezza;
public:
    Rettangolo() { base = 1; altezza = 1; }                  // costruttore di default
    Rettangolo(double b, double h) { base = b; altezza = h; }
};
Rettangolo r1;            // chiama Rettangolo()
Rettangolo r2(3, 4);      // chiama Rettangolo(double, double)
Rettangolo* p = new Rettangolo(5, 2);   // anche new chiama il costruttore
\end{lstlisting}

\textbf{Il costruttore di default} è quello che si può chiamare \textbf{senza
argomenti}. Se nella classe non scrivi \emph{nessun} costruttore, il compilatore
ne genera uno vuoto, che \emph{non} inizializza gli attributi di tipo base
(contengono valori casuali). Se invece scrivi almeno un costruttore con
parametri, quello automatico \textbf{sparisce}, e \cc|Rettangolo r;| oppure
\cc|Rettangolo v[10];| smettono di compilare. È la trappola del Modulo 1 sugli
array di oggetti.

\textbf{Parametri di default.} Un solo costruttore può coprire più casi dando un
valore predefinito ai parametri finali:
\cc|Rettangolo(double b = 1, double h = 1)| si può chiamare come
\cc|Rettangolo()|, \cc|Rettangolo(4)| o \cc|Rettangolo(3, 4)|. I default vanno
scritti da destra a sinistra: non puoi dare un default al primo parametro e non
al secondo.

\subsection*{La lista di inizializzazione}

Dopo i parametri del costruttore, prima del corpo, si può scrivere una
\textbf{lista di inizializzazione}, introdotta dai due punti:
\begin{lstlisting}
Rettangolo(double b, double h) : base(b), altezza(h) {
    // corpo (qui gli attributi hanno già i loro valori)
}
\end{lstlisting}
La differenza con l'assegnare nel corpo è sottile ma importante. Gli attributi
vengono \textbf{creati} prima che inizi il corpo del costruttore: con la lista li
crei direttamente con il valore giusto; con le assegnazioni nel corpo prima li crei
``vuoti'', poi ci scrivi sopra. Per un \cc|int| non cambia quasi nulla, ma ci sono
tre casi in cui la lista è \textbf{obbligatoria}:
\begin{itemize}
  \item attributi \cc|const| (non si possono assegnare dopo la nascita);
  \item attributi reference (vanno legati alla nascita, Modulo 3);
  \item attributi che sono \textbf{oggetti} di una classe senza costruttore di
        default (vedi l'Esercizio~4.5 con \cc|Punto2D| dentro \cc|Rettangolo|).
\end{itemize}
Gli attributi vengono inizializzati nell'\textbf{ordine in cui sono dichiarati
nella classe}, non nell'ordine in cui li scrivi nella lista. Conviene quindi
scriverli nella lista nello stesso ordine, per non confondersi.

Altri due modi di inizializzare, presenti nelle slide del corso:
\begin{lstlisting}
class MyClass {
    int x = 18;            // inizializzazione nella dichiarazione (valore di partenza)
    char c {'M'};          // inizializzazione uniforme con le graffe
public:
    MyClass(int a, char b) : x(a), c(b) {}
    MyClass(char b) : MyClass(10, b) {}     // costruttore con DELEGA: ne riusa un altro
};
\end{lstlisting}

\subsection*{Il distruttore}

\begin{definizione}[: distruttore]
Il \textbf{distruttore} è un metodo speciale eseguito \textbf{automaticamente}
quando un oggetto muore. Si chiama come la classe preceduta da una tilde,
\cc|~Rettangolo()|, non ha tipo di ritorno, \textbf{non ha parametri} e ce n'è
\textbf{uno solo} per classe. Se non lo scrivi, il compilatore ne crea uno vuoto.
\end{definizione}

A cosa serve? Per un oggetto che contiene solo numeri e stringhe, a niente: quello
generato dal compilatore va benissimo. Diventa \textbf{indispensabile} quando
l'oggetto \textbf{possiede risorse}: memoria presa con \cc|new|, un file aperto,
una connessione. La regola è: \textbf{quello che il costruttore prende, il
distruttore restituisce}.
\begin{lstlisting}
class Vettore {
    int* dati;
    int n;
public:
    Vettore(int dim) : dati(new int[dim]), n(dim) {}   // prende memoria
    ~Vettore() { delete[] dati; }                      // la restituisce
};
\end{lstlisting}
Il grande vantaggio: chi usa \cc|Vettore| non deve mai ricordarsi di fare
\cc|delete[]|. Quando l'oggetto esce dal suo blocco, il distruttore parte da solo,
anche se si esce per un'eccezione. Questa idea (legare una risorsa alla vita di un
oggetto) si chiama \textbf{RAII} ed è il modo C++ di evitare i memory leak.

\begin{analogia}[: la camera d'albergo]
Il costruttore è il check-in: ricevi la chiave della camera (la memoria). Il
distruttore è il check-out: restituisci la chiave. Un oggetto senza distruttore,
che però ha preso una camera con \cc|new|, è un cliente che se ne va con la chiave
in tasca: la camera resta occupata per sempre (memory leak).
\end{analogia}

\subsection*{Il puntatore this}

Il codice di un metodo esiste una volta sola, condiviso da tutti gli oggetti
(Modulo 1). Allora, quando scrivi \cc|a.setX(3)|, come fa \cc|setX| a sapere di
dover modificare la \cc|x| di \cc|a| e non quella di \cc|b|? Il compilatore passa
al metodo, di nascosto, un parametro in più: l'\textbf{indirizzo dell'oggetto} su
cui è stato chiamato. Dentro il metodo questo indirizzo si chiama \cc|this|.

\begin{definizione}[: this]
\cc|this| è un puntatore, disponibile dentro ogni metodo non statico, che contiene
l'indirizzo dell'oggetto su cui il metodo è stato invocato. \cc|*this| è l'oggetto
stesso.
\end{definizione}

Usi tipici:
\begin{itemize}
  \item \textbf{distinguere} un attributo da un parametro con lo stesso nome:
        \cc|this->x = x;| (a sinistra l'attributo, a destra il parametro);
  \item \textbf{restituire l'oggetto stesso} per concatenare le chiamate:
        \cc|Punto& setX(int x) { this->x = x; return *this; }| permette
        \cc|p.setX(3).setY(4);|. È lo stesso meccanismo di \cc|cout << a << b|;
  \item \textbf{riconoscere se stesso}: \cc|if (this == &altro)| (stesso
        indirizzo, quindi stesso oggetto). È indispensabile nell'operatore di
        assegnamento, come vedremo.
\end{itemize}

\subsection*{Metodi e oggetti const}

Lo abbiamo anticipato nel Modulo 1: un metodo \cc|const| promette di non
modificare l'oggetto. Ora si capisce meglio il perché: dentro un metodo \cc|const|,
\cc|this| è un puntatore a oggetto \emph{costante}, e il compilatore blocca ogni
scrittura sugli attributi. Su un oggetto dichiarato \cc|const|
(\cc|const Rettangolo r(2, 2);|) o ricevuto come \cc|const Rettangolo&| si possono
chiamare \textbf{solo} i metodi \cc|const|.

\subsection*{Membri static}

Normalmente ogni oggetto ha la sua copia di ogni attributo. A volte però serve un
dato che appartiene \textbf{alla classe intera} e non al singolo oggetto: quanti
oggetti esistono, un contatore di codici progressivi, una costante comune.

\begin{definizione}[: membro statico]
Un attributo \cc|static| esiste in \textbf{una sola copia} condivisa da tutti gli
oggetti della classe, ed esiste anche se non è stato creato nessun oggetto. Un
metodo \cc|static| si può chiamare senza oggetto (\cc|Classe::metodo()|) e
quindi \textbf{non ha \texttt{this}}: può usare solo gli attributi statici.
\end{definizione}

\begin{lstlisting}
class Punto {
    int x, y;
    static int vivi;                        // dichiarazione (dentro la classe)
public:
    Punto() : x(0), y(0) { vivi++; }
    ~Punto() { vivi--; }
    static int getVivi() { return vivi; }
};
int Punto::vivi = 0;                        // definizione (fuori, una sola volta)
cout << Punto::getVivi();                   // si chiama sulla classe
\end{lstlisting}
L'attributo statico va \textbf{definito} fuori dalla classe, in un \texttt{.cpp}:
dimenticarlo dà un errore del \textbf{linker} (\emph{undefined reference to
`Punto::vivi'}). Fanno eccezione le costanti intere (\cc|static const int MAX = 10;|),
che si possono inizializzare direttamente nella classe.

\subsection*{Copiare un oggetto: il problema}

Gli oggetti si possono copiare, come le \cc|struct| del C: \cc|Vettore b = a;|
oppure \cc|b = a;|. Se non scrivi nulla, il compilatore copia \textbf{membro a
membro} (le slide dicono ``bit a bit''): ogni attributo di \cc|a| viene copiato
nel corrispondente attributo di \cc|b|. Per numeri e \cc|string| è perfetto.
Ma se un attributo è un \textbf{puntatore}, viene copiato il puntatore, cioè
\textbf{l'indirizzo}, non il blocco a cui punta. Si ottengono così due oggetti che
condividono lo stesso blocco di heap: è la \textbf{copia superficiale}
(\emph{shallow copy}).

Perché è un disastro?
\begin{enumerate}
  \item \textbf{Modifiche incrociate}: cambi un elemento di \cc|b| e cambia anche
        \cc|a|, perché è lo stesso blocco (Esercizio~4.7).
  \item \textbf{Doppio delete}: quando muoiono, \emph{entrambi} i distruttori fanno
        \cc|delete[]| dello \emph{stesso} blocco. Il secondo \cc|delete[]| è
        comportamento indefinito: spesso il programma va in crash alla fine.
  \item \textbf{Dangling pointer}: se \cc|a| muore prima di \cc|b|, \cc|b| si
        ritrova con un puntatore a memoria già liberata.
\end{enumerate}
La soluzione è la \textbf{copia profonda} (\emph{deep copy}): la copia si prende
un blocco \textbf{tutto suo} e ci copia dentro gli elementi uno per uno. Per
ottenerla bisogna scrivere due metodi speciali.

\begin{analogia}[: la chiave e la casa]
Copia superficiale = fare un duplicato della \textbf{chiave}: ora due persone
entrano nella stessa casa. Se uno sposta i mobili, l'altro li trova spostati; se
uno vende la casa, l'altro ha in mano una chiave che apre la casa di uno
sconosciuto. Copia profonda = costruire una \textbf{casa identica} e dare
all'altro la chiave della casa nuova.
\end{analogia}

\subsection*{Il costruttore di copia}

\begin{definizione}[: costruttore di copia]
È il costruttore che crea un oggetto \textbf{nuovo} come copia di uno esistente.
Ha la forma \cc|Classe(const Classe& altro)|.
\end{definizione}
Viene chiamato automaticamente in tre situazioni:
\begin{itemize}
  \item quando un oggetto nasce da un altro: \cc|Vettore b = a;| oppure \cc|Vettore b(a);|
        (attenzione: anche se c'è il simbolo \cc|=|, qui \cc|b| sta \emph{nascendo},
        quindi è un costruttore e non un assegnamento);
  \item quando un oggetto viene \textbf{passato per valore} a una funzione (il
        parametro è una copia);
  \item quando una funzione \textbf{restituisce un oggetto per valore} (il
        compilatore spesso ottimizza ed evita la copia, ma la regola è questa).
\end{itemize}
Perché il parametro è una \textbf{reference}? Se fosse \cc|Classe(Classe altro)|,
per passare \cc|altro| per valore servirebbe... una copia, cioè di nuovo il
costruttore di copia, che richiederebbe un'altra copia, all'infinito. Il
compilatore lo vieta. Perché \cc|const|? Perché copiare non deve modificare
l'originale, e così si possono copiare anche oggetti costanti.

\subsection*{L'operatore di assegnamento}

\cc|b = a;| quando \cc|b| \textbf{esiste già} non chiama il costruttore di copia,
ma l'\textbf{operatore di assegnamento} \cc|operator=|. La differenza è
fondamentale: \cc|b| possiede già un suo blocco di heap, che va
\textbf{restituito} prima di prenderne uno nuovo, altrimenti si perde (leak).

\begin{passopasso}[: scrivere operator=]
\begin{lstlisting}
Vettore& operator=(const Vettore& altro) {
    if (this == &altro) return *this;   // 1. a = a: non fare nulla
    delete[] dati;                      // 2. libera il TUO vecchio blocco
    n = altro.n;                        // 3. copia gli attributi semplici
    dati = new int[n];                  // 4. prendi un blocco nuovo
    for (int i = 0; i < n; i++)         // 5. copia gli elementi
        dati[i] = altro.dati[i];
    return *this;                       // 6. restituisci te stesso
}
\end{lstlisting}
\begin{enumerate}
  \item Il controllo di \textbf{autoassegnamento} serve perché, se \cc|this| e
        \cc|altro| sono lo stesso oggetto, il passo 2 distruggerebbe proprio i dati
        da copiare.
  \item Si restituisce \cc|*this| per reference (\cc|Vettore&|) così funziona
        \cc|a = b = c;|, che si legge \cc|a = (b = c);|.
  \item Variante più sicura (usata nel laboratorio): prima si alloca e si copia
        in un blocco nuovo, \emph{poi} si libera il vecchio. Così, se \cc|new|
        fallisse, l'oggetto resterebbe intatto.
\end{enumerate}
\end{passopasso}

\subsection*{La Regola dei Tre}

\begin{definizione}[: Regola dei Tre]
Se una classe ha bisogno di definire \textbf{uno} tra distruttore, costruttore di
copia e operatore di assegnamento, quasi sicuramente ha bisogno di definirli
\textbf{tutti e tre}.
\end{definizione}
Il motivo è che tutti e tre nascono dallo stesso fatto: la classe possiede una
risorsa (di solito memoria dinamica). Se serve il distruttore per liberarla, le
versioni automatiche di copia e assegnamento la condividerebbero (copia
superficiale), e si torna ai disastri visti sopra. All'esame, ogni volta che vedi
un attributo puntatore con \cc|new| nel costruttore, scrivi subito tutti e tre.

\subsection*{Funzioni friend (anteprima)}

Una funzione esterna dichiarata \cc|friend| dentro una classe può accedere ai
membri privati di quella classe. Sembra una violazione dell'information hiding,
ma non lo è: è la classe stessa a decidere chi sono i suoi ``amici'', nella sua
dichiarazione. Si usa soprattutto per gli operatori \cc|<<|, \cc|>>|, \cc|==| e
\cc|+|, che studieremo nel Modulo 10.

\subsection*{Le eccezioni}

Finora, quando un metodo riceveva dati sbagliati, li abbiamo corretti o abbiamo
restituito \cc|false|. Non sempre si può: un costruttore non restituisce niente, e
\cc|int get(int i)| non ha un valore ``speciale'' per dire ``indice sbagliato''.
Fermare il programma con \cc|exit(1)| è brutale. Il C++ offre un meccanismo
apposito: le \textbf{eccezioni}.

\begin{itemize}
  \item \cc|throw valore;| \textbf{lancia} un'eccezione: la funzione si interrompe
        all'istante, e così tutte le funzioni chiamanti, finché qualcuno non la cattura.
  \item \cc|try { ... }| delimita il codice ``sorvegliato'';
  \item \cc|catch (const Tipo& e) { ... }| subito dopo il \cc|try| \textbf{cattura}
        le eccezioni di quel tipo e le gestisce. Possono esserci più \cc|catch|,
        provati in ordine; \cc|catch (...)| cattura qualunque cosa.
\end{itemize}
\begin{lstlisting}
#include <stdexcept>
void setOre(int h) {
    if (h < 0 || h > 23) throw invalid_argument("ore valide: 0-23");
    ore = h;
}
try {
    t.setOre(100);                 // lancia
    cout << "mai eseguito";        // saltato
} catch (const invalid_argument& e) {
    cout << "Errore: " << e.what(); // what() restituisce il messaggio
}
\end{lstlisting}
La libreria \cc|<stdexcept>| fornisce eccezioni pronte: \cc|invalid_argument|
(argomento non valido), \cc|out_of_range| (indice fuori dai limiti),
\cc|underflow_error|, \cc|overflow_error|, \cc|runtime_error|. Sono tutte
``figlie'' di \cc|exception| (Modulo 9): \cc|catch (const exception& e)| le prende
tutte. Si catturano per \textbf{reference costante}, come consigliano le slide.

Se un'eccezione esce da un \textbf{costruttore}, l'oggetto \textbf{non viene
creato}: è il modo pulito per impedire la nascita di oggetti non validi.

\begin{inparole}
\begin{itemize}
  \item Costruttore = collaudo (nasce l'oggetto); distruttore = check-out (muore).
        Partono da soli.
  \item Se scrivi un costruttore con parametri, quello senza parametri non esiste più.
  \item Lista di inizializzazione \cc|: a(x), b(y)|: obbligatoria per \cc|const|,
        reference e oggetti senza costruttore di default.
  \item \cc|this| = indirizzo dell'oggetto su cui stai lavorando.
  \item \cc|static| = uno solo per tutta la classe; va definito fuori.
  \item Puntatore come attributo + \cc|new| nel costruttore $\Rightarrow$
        \textbf{Regola dei Tre}: distruttore, costruttore di copia, \cc|operator=|.
  \item \cc|Classe b = a;| $\to$ costruttore di copia; \cc|b = a;| con \cc|b| già
        esistente $\to$ \cc|operator=|.
  \item Errori che non si possono gestire lì: \cc|throw|; chi può gestirli:
        \cc|try|/\cc|catch|.
\end{itemize}
\end{inparole}

% ---------------------------------------------------------------------
\section{Rappresentazione Visiva della Memoria (TikZ)}
% ---------------------------------------------------------------------

\textbf{Nascita e morte degli oggetti locali} (Esercizio~4.2): lo stack cresce
quando gli oggetti nascono e si svuota in ordine inverso.
\begin{center}
\begin{tikzpicture}
  \foreach \nome/\y in {A/0,B/0.8,C/1.6}
    \node[cella,fill=stackcol,minimum width=2.2cm] (o\nome) at (0,\y) {\nome};
  \node[etichetta,below=0.15cm of oA] {stack del \texttt{main}};
  \draw[freccia] (-2.4,-0.2) -- node[left,etichetta,text width=1.7cm,align=right]{nascono\\ A, B, C} (-2.4,1.8);
  \draw[freccia] (2.4,1.8) -- node[right,etichetta,text width=2.6cm]{muoiono\\ C, B (fine blocco)\\ poi A (fine \texttt{main})} (2.4,-0.2);
\end{tikzpicture}
\end{center}

\needspace{7cm}
\textbf{Copia superficiale contro copia profonda.} A sinistra quello che fa il
compilatore da solo, a destra quello che fa il costruttore di copia scritto da noi.
\begin{center}
\begin{tikzpicture}
  % superficiale
  \node[font=\bfseries\small] at (2.2,2.4) {SUPERFICIALE (sbagliata)};
  \node[cella,fill=stackcol] (a1) at (0,1.2) {a.dati};
  \node[cella,fill=stackcol] (b1) at (0,-0.2) {b.dati};
  \foreach \v [count=\i] in {10,20,30}
    \node[cellina,fill=heapcol] (h\i) at (2.6+\i*0.75,0.5) {\v};
  \draw[freccia] (a1.east) -- (h1.north west);
  \draw[frecciarossa] (b1.east) -- (h1.south west);
  \node[etichetta,text=red!75!black,below=0.25cm of h2,text width=4cm,align=center]
       {stesso blocco: modifiche incrociate e doppio \texttt{delete[]}};
  % profonda
  \node[font=\bfseries\small] at (10.2,2.4) {PROFONDA (giusta)};
  \node[cella,fill=stackcol] (a2) at (8,1.2) {a.dati};
  \node[cella,fill=stackcol] (b2) at (8,-0.2) {b.dati};
  \foreach \v [count=\i] in {10,20,30}
    \node[cellina,fill=heapcol] (p\i) at (10.6+\i*0.75,1.2) {\v};
  \foreach \v [count=\i] in {10,20,30}
    \node[cellina,fill=heapcol] (q\i) at (10.6+\i*0.75,-0.2) {\v};
  \draw[freccia] (a2.east) -- (p1.west);
  \draw[freccia] (b2.east) -- (q1.west);
  \node[etichetta,below=0.25cm of q2,text width=4cm,align=center]
       {ognuno il suo blocco: indipendenti};
\end{tikzpicture}
\end{center}

\needspace{6cm}
\textbf{Cosa fa \texttt{c = a;}} con \cc|operator=| scritto bene (\cc|c| aveva 5
elementi, \cc|a| ne ha 3).
\begin{center}
\begin{tikzpicture}
  \node[cella,fill=stackcol] (c) at (0,0) {c.dati};
  \foreach \i in {1,...,5}
    \node[cellina,fill=heapcol] (v\i) at (2+\i*0.75,0.9) {0};
  \draw[red!75!black,thick] ($(v1.south west)+(-0.05,-0.05)$) -- ($(v5.north east)+(0.05,0.05)$);
  \foreach \v [count=\i] in {10,20,30}
    \node[cellina,fill=heapcol] (w\i) at (2+\i*0.75,-0.9) {\v};
  \draw[freccia,dashed,black!50] (c.east) -- (v1.west);
  \draw[freccia] (c.east) -- (w1.west);
  \node[text width=6.6cm,font=\small,anchor=west] at (7.2,0) {%
    \textbf{1.} \texttt{this != \&a}? sì, si procede\\
    \textbf{2.} \texttt{delete[] dati} (vecchio blocco da 5, barrato)\\
    \textbf{3--5.} \texttt{n = 3}, nuovo blocco da 3, copia\\
    \textbf{6.} \texttt{return *this}};
\end{tikzpicture}
\end{center}

\needspace{6cm}
\textbf{Un attributo static è uno solo.} Ogni oggetto ha i suoi \cc|x| e \cc|y|;
\cc|vivi| sta fuori dagli oggetti, nella memoria statica, condiviso da tutti.
\begin{center}
\begin{tikzpicture}
  \foreach \nome/\xx/\yy [count=\i] in {a/1/2,b/0/0,c/1/2}{
    \node[cella,fill=privcol,minimum width=1.6cm] (x\i) at (\i*3,0.4) {x = \xx};
    \node[cella,fill=privcol,minimum width=1.6cm,below=0cm of x\i] (y\i) {y = \yy};
    \node[font=\small\bfseries,above=0cm of x\i] {\nome};
  }
  \node[cella,fill=green!15!sfondo,minimum width=2.4cm] (vivi) at (6,-2.3) {Punto::vivi = 3};
  \node[etichetta,below=0.05cm of vivi] {memoria statica: una sola copia};
  \foreach \i in {1,2,3} \draw[freccia,dashed] (y\i.south) -- (vivi.north);
\end{tikzpicture}
\end{center}

% ---------------------------------------------------------------------
\section{Pattern di Codice di Laboratorio}
% ---------------------------------------------------------------------

La classe \cc|Vettore| è lo scheletro di \textbf{ogni} classe che possiede memoria
dinamica: costruttore che alloca, distruttore che libera, costruttore di copia e
\cc|operator=| profondi, accesso controllato con eccezioni. I messaggi tra
parentesi quadre mostrano quale metodo speciale viene chiamato e quando.

\codice{codice/m04/pattern.cpp}
\uscita{codice/m04/output/pattern.txt}

\textbf{Cosa osservare:}
\begin{itemize}
  \item \cc|Vettore b = a;| stampa \texttt{[costruttore di copia]}, non
        \texttt{[operator=]}: \cc|b| sta nascendo.
  \item Dopo \cc|b.set(0, 99)|, \cc|a| è rimasto \texttt{10 20 30}: la copia è
        profonda.
  \item \cc|c = a;| stampa \texttt{[operator=]}: \cc|c| esisteva già, con 5 elementi.
  \item I tre distruttori partono da soli alla fine del \cc|main|, in ordine inverso
        di creazione, e ciascuno libera il \emph{suo} blocco: nessun doppio \cc|delete|.
\end{itemize}

% ---------------------------------------------------------------------
\section{Tips \& Trappole da Esame}
% ---------------------------------------------------------------------

\begin{trappola}[: costruttore di default sparito]
Appena scrivi un costruttore con parametri, il compilatore non genera più quello
senza parametri. \cc|Classe x;|, \cc|Classe v[10];| e \cc|new Classe[n]| smettono
di compilare (\emph{no matching function for call to 'Classe::Classe()'}).
Soluzione: aggiungi \cc|Classe()| o dai valori di default ai parametri.
\end{trappola}

\begin{trappola}[: la copia superficiale]
Classe con un attributo puntatore e \cc|new| nel costruttore, ma senza costruttore
di copia e \cc|operator=|: ogni copia condivide il blocco. Il sintomo tipico è un
crash \textbf{alla fine del programma} (doppio \cc|delete|) o dopo aver passato
l'oggetto \textbf{per valore} a una funzione. Regola dei Tre, sempre.
\end{trappola}

\begin{trappola}[: inizializzazione o assegnamento?]
\cc|Classe b = a;| è una \textbf{inizializzazione} (costruttore di copia), anche se
c'è un \cc|=|. \cc|Classe b; b = a;| è un \textbf{assegnamento} (\cc|operator=|).
Domanda classicissima nei quiz.
\end{trappola}

\begin{trappola}[: operator= senza controllo e senza delete]
Senza \cc|if (this == &altro)|, l'istruzione \cc|a = a;| libera i dati e poi prova a
copiarli dal blocco appena liberato. Senza \cc|delete[]| del vecchio blocco, ogni
assegnamento produce un memory leak.
\end{trappola}

\begin{trappola}[: parametro del costruttore di copia per valore]
\cc|Classe(Classe altro)| non compila: per creare \cc|altro| servirebbe il
costruttore di copia stesso (ricorsione infinita). Il parametro è sempre
\cc|const Classe&|.
\end{trappola}

\begin{trappola}[: attributo static non definito]
\cc|static int vivi;| dentro la classe è solo una \emph{dichiarazione}. Senza
\cc|int Punto::vivi = 0;| fuori dalla classe, il linker dà \emph{undefined reference}.
\end{trappola}

\begin{trappola}[: nomi già usati dal sistema]
Su MinGW un metodo chiamato \cc|alloca()| non compila con un errore
incomprensibile (\emph{expected identifier before '(' token}): \cc|alloca| è già il
nome di una funzione di sistema, definita come macro. È successo scrivendo
l'Esercizio~4.9. Se un errore non ha senso su un nome semplice, prova a rinominarlo.
\end{trappola}

\begin{consiglio}[: stampare nei costruttori e nei distruttori]
Quando non capisci cosa succede, aggiungi temporaneamente un \cc|cout| in
costruttore, costruttore di copia, \cc|operator=| e distruttore, come nel pattern:
vedrai esattamente quando nasce, si copia e muore ogni oggetto.
\end{consiglio}

\begin{consiglio}[: passare gli oggetti per const reference]
\cc|void f(Vettore v)| fa una copia profonda a ogni chiamata (lento, e un disastro se
manca il costruttore di copia). \cc|void f(const Vettore& v)| non copia niente.
Per valore solo se la funzione ha davvero bisogno di una copia da modificare.
\end{consiglio}

% ---------------------------------------------------------------------
\section{Quiz a Crocette (Prova Scritta)}
% ---------------------------------------------------------------------

\begin{quiz}{Quale affermazione sul costruttore è \textbf{falsa}?}
  \item ha lo stesso nome della classe
  \item restituisce \texttt{void}
  \item può essere sovraccaricato (overloading)
  \item viene eseguito automaticamente alla creazione dell'oggetto
\end{quiz}

\begin{quiz}{Una classe ha solo il costruttore \texttt{A(int x)}. Quale istruzione \textbf{non} compila?}
  \item \texttt{A a(5);}
  \item \texttt{A* p = new A(3);}
  \item \texttt{A v[10];}
  \item \texttt{A b = A(2);}
\end{quiz}

\begin{quiz}{Quante volte e con quanti parametri si può definire il distruttore di una classe?}
  \item una volta sola, senza parametri
  \item una volta per ogni costruttore
  \item più volte, purché con parametri diversi
  \item una volta sola, con un parametro di tipo \texttt{int}
\end{quiz}

\begin{quiz}{Dato \texttt{Vettore a(3);} quale istruzione chiama il \textbf{costruttore di copia}?}
  \item \texttt{Vettore b; b = a;}
  \item \texttt{Vettore b = a;}
  \item \texttt{a = a;}
  \item \texttt{Vettore* p = \&a;}
\end{quiz}

\begin{quiz}{Perché il parametro del costruttore di copia è \texttt{const Classe\&} e non \texttt{Classe}?}
  \item per motivi di velocità, ma funzionerebbe anche per valore
  \item perché il passaggio per valore richiederebbe a sua volta una copia, causando una ricorsione infinita
  \item perché i costruttori non possono ricevere oggetti
  \item perché le reference sono obbligatorie in tutti i costruttori
\end{quiz}

\begin{quiz}{Cosa contiene \texttt{this} dentro un metodo?}
  \item una copia dell'oggetto su cui il metodo è stato chiamato
  \item l'indirizzo dell'oggetto su cui il metodo è stato chiamato
  \item l'indirizzo della classe
  \item l'indirizzo del metodo stesso
\end{quiz}

\begin{quiz}{Un attributo dichiarato \texttt{static} in una classe:}
  \item è condiviso da tutti gli oggetti della classe
  \item non può essere modificato
  \item viene allocato nello stack di ogni oggetto
  \item è accessibile solo dai metodi \texttt{const}
\end{quiz}

\begin{quiz}{Cosa afferma la Regola dei Tre?}
  \item una classe deve avere al massimo tre costruttori
  \item se serve uno tra distruttore, costruttore di copia e operatore di assegnamento, servono tutti e tre
  \item ogni classe deve avere almeno tre metodi pubblici
  \item un oggetto può essere copiato al massimo tre volte
\end{quiz}

\begin{quiz}{Una classe ha un attributo \texttt{int* dati} allocato con \texttt{new[]} nel costruttore e un distruttore che fa \texttt{delete[]}, ma \textbf{nessun} costruttore di copia. Cosa succede con \texttt{\{ Classe a; Classe b = a; \}}?}
  \item niente di strano: il compilatore fa una copia profonda
  \item \texttt{a} e \texttt{b} condividono il blocco e alla fine il blocco viene liberato due volte
  \item errore di compilazione
  \item \texttt{b} riceve un puntatore nullo
\end{quiz}

\begin{quiz}{Perché \texttt{operator=} controlla \texttt{if (this == \&altro)}?}
  \item per evitare di copiare oggetti di classi diverse
  \item perché nell'autoassegnamento il \texttt{delete} iniziale distruggerebbe i dati da copiare
  \item per velocizzare la copia di oggetti uguali
  \item perché il compilatore lo richiede
\end{quiz}

\begin{quiz}{Un costruttore lancia un'eccezione con \texttt{throw}. Cosa succede all'oggetto?}
  \item viene creato con gli attributi a zero
  \item viene creato normalmente, l'eccezione è solo un avviso
  \item non viene creato
  \item il programma termina sempre
\end{quiz}

\begin{soluzione}[: risposte commentate]
\begin{enumerate}[label=\textbf{\arabic*.}]
  \item \textbf{B}. Il costruttore non ha \emph{nessun} tipo di ritorno, nemmeno \cc|void|.
  \item \textbf{C}. L'array richiede il costruttore senza parametri, che non esiste più.
  \item \textbf{A}. Un solo distruttore, senza parametri e senza tipo di ritorno.
  \item \textbf{B}. \cc|b| nasce da \cc|a|. La A e la C sono assegnamenti
        (\cc|operator=|), la D copia solo un indirizzo.
  \item \textbf{B}. È proprio un divieto, non una questione di velocità.
  \item \textbf{B}. \cc|this| è un puntatore; \cc|*this| è l'oggetto.
  \item \textbf{A}. Esiste in un'unica copia, nella memoria statica, anche senza oggetti.
  \item \textbf{B}. I tre metodi gestiscono la stessa risorsa: o servono tutti o nessuno.
  \item \textbf{B}. La copia automatica è superficiale: doppio \cc|delete[]| alla
        chiusura del blocco. Compila senza problemi (C è falsa).
  \item \textbf{B}. \cc|a = a|: \cc|delete[] dati| cancellerebbe proprio i dati di \cc|altro|.
  \item \textbf{C}. Un oggetto la cui costruzione fallisce non esiste: è il modo pulito
        per impedire la nascita di oggetti non validi. Il programma termina solo se
        nessuno cattura l'eccezione.
\end{enumerate}
\end{soluzione}

% ---------------------------------------------------------------------
\section{Laboratorio a Step (Simulazione d'Esame)}
% ---------------------------------------------------------------------

\begin{traccia}[: playlist musicale]
Si vuole gestire una raccolta di playlist. Ogni playlist ha un nome e contiene un
numero variabile di brani, memorizzati in un array dinamico (niente STL). Ogni brano
ha titolo, artista e durata in secondi.
\end{traccia}

\begin{step}{la classe Brano}
Scrivi \cc|Brano| con titolo, artista e durata. Prevedi: un costruttore di default
(\q{?}, \q{?}, 0), un costruttore completo con \textbf{lista di inizializzazione}
(durata negativa $\to$ 0) e un costruttore con solo il titolo che
\textbf{delega} al completo con artista \q{Sconosciuto} e durata 180. Getter
\cc|const| e \cc|stampa()| nel formato \texttt{Titolo - Artista (m:ss)}.\\[0.3em]
\textbf{Concetto nuovo:} overloading dei costruttori, lista di inizializzazione, delega.
\end{step}

\begin{step}{la classe Playlist}
\cc|Playlist| ha un nome e un array dinamico di \cc|Brano| (\cc|Brano* brani|), con
\cc|n| brani presenti e capacità \cc|cap|. Il costruttore riceve il nome e alloca
spazio per 2 brani; il \textbf{distruttore} libera l'array. Fai stampare a entrambi
un messaggio.\\[0.3em]
\textbf{Concetto nuovo:} distruttore, risorsa acquisita nel costruttore e
rilasciata nel distruttore.
\end{step}

\begin{step}{aggiungere brani}
Aggiungi \cc|void aggiungi(const Brano& b)|, che raddoppia la capacità quando
serve tramite un metodo \textbf{privato} \cc|cresci()|, e \cc|int durataTotale() const|.\\[0.3em]
\textbf{Concetto nuovo:} riallocazione dentro una classe, parametro \cc|const&|.
\end{step}

\begin{step}{copiare una playlist}
Scrivi il \textbf{costruttore di copia} profondo: la copia ha come nome quello
originale seguito da \q{ (copia)} e un array tutto suo. Nel \cc|main| crea una
copia dentro un blocco \cc|{ }|, aggiungile un brano e verifica che l'originale
non cambi e che la copia venga distrutta alla fine del blocco.\\[0.3em]
\textbf{Concetto nuovo:} costruttore di copia, copia profonda.
\end{step}

\begin{step}{assegnare una playlist}
Scrivi \cc|operator=| (che copia i brani ma \textbf{non} il nome), con controllo di
autoassegnamento. Prova \cc|r = p;| e \cc|r = r;|.\\[0.3em]
\textbf{Concetto nuovo:} operatore di assegnamento, \cc|this|, Regola dei Tre.
\end{step}

\begin{step}{contatore ed eccezioni}
Aggiungi un attributo \cc|static int vive| che conta le playlist esistenti
(aggiornato da costruttori e distruttore) e un metodo statico \cc|quanteVive()|.
Aggiungi \cc|Brano& at(int i)|, che lancia \cc|out_of_range| se l'indice non è
valido, e gestisci l'errore nel \cc|main| con \cc|try|/\cc|catch|.\\[0.3em]
\textbf{Concetto nuovo:} membri \cc|static|, \cc|throw|, \cc|try|, \cc|catch|.
\end{step}

\subsection*{Soluzione completa}
\codice{codice/m04/lab.cpp}
\uscita{codice/m04/output/lab.txt}

\begin{consiglio}[: cosa notare nella soluzione]
\begin{itemize}
  \item \cc|new Brano[cap]| crea \cc|cap| brani con il costruttore di
        \textbf{default}: senza \cc|Brano()| non compilerebbe.
  \item Nello step 4 il messaggio \texttt{distrutta \q{Classici (copia)}} compare
        \emph{alla chiusura del blocco}, prima dello step 5: è il distruttore che
        parte da solo.
  \item Alla fine del \cc|main| il contatore torna a 0: ogni costruttore (anche
        quello di copia) incrementa, ogni distruttore decrementa.
  \item In \cc|operator=| si alloca \emph{prima} e si libera \emph{dopo}: se la
        copia fallisse, la playlist resterebbe intatta.
  \item \cc|to_string(i)| (da \cc|<string>|) converte un numero in stringa per il
        messaggio dell'eccezione.
\end{itemize}
\end{consiglio}

% ---------------------------------------------------------------------
\section{Domande e Risposte Modello per l'Orale}
% ---------------------------------------------------------------------

\begin{orale}[D1. Cosa sono costruttore e distruttore?]
Sono metodi speciali invocati automaticamente. Il costruttore viene eseguito
quando un oggetto viene creato: ha lo stesso nome della classe, nessun tipo di
ritorno, può avere parametri e può essere sovraccaricato; serve a portare
l'oggetto in uno stato iniziale valido. Il distruttore viene eseguito quando
l'oggetto viene distrutto (fine del blocco, \cc|delete|, fine del programma): si
chiama come la classe preceduto da \texttt{\textasciitilde}, non ha parametri né
tipo di ritorno ed è unico. Serve a rilasciare le risorse acquisite dall'oggetto,
tipicamente la memoria dinamica allocata nel costruttore.
\end{orale}

\begin{orale}[D2. Differenza tra copia superficiale e copia profonda?]
Se una classe non definisce il costruttore di copia, il compilatore copia gli
attributi membro a membro. Se un attributo è un puntatore viene copiato
l'indirizzo, e i due oggetti condividono lo stesso blocco di heap: è la copia
superficiale. Le conseguenze sono modifiche incrociate, dangling pointer e doppia
deallocazione da parte dei distruttori. La copia profonda alloca invece per la
copia un nuovo blocco e vi copia il contenuto, rendendo i due oggetti indipendenti;
si ottiene definendo il costruttore di copia e l'operatore di assegnamento.
\end{orale}

\begin{orale}[D3. Quando viene chiamato il costruttore di copia e quando operator=?]
Il costruttore di copia viene chiamato quando un oggetto \textbf{nasce} come copia di
un altro: inizializzazione (\cc|Classe b = a;| o \cc|Classe b(a);|), passaggio di
un oggetto per valore e restituzione per valore. L'operatore di assegnamento viene
chiamato quando a un oggetto \textbf{già esistente} se ne assegna un altro
(\cc|b = a;|). Per questo \cc|operator=| deve prima liberare le risorse che
l'oggetto possiede già e deve gestire l'autoassegnamento.
\end{orale}

\begin{orale}[D4. Cos'è la Regola dei Tre?]
Afferma che se una classe ha bisogno di definire esplicitamente uno tra
distruttore, costruttore di copia e operatore di assegnamento, allora deve
definirli tutti e tre. Il motivo è che la necessità nasce dalla gestione di una
risorsa, tipicamente memoria dinamica: il distruttore la rilascia, e le versioni
generate dal compilatore di copia e assegnamento la condividerebbero, causando
copie superficiali e doppie deallocazioni.
\end{orale}

\begin{orale}[D5. Cos'è il puntatore this e quando si usa?]
\cc|this| è un puntatore implicito, disponibile in ogni metodo non statico, che
contiene l'indirizzo dell'oggetto su cui il metodo è stato invocato; è il modo in
cui un unico codice dei metodi lavora sui dati dell'oggetto giusto. Si usa per
distinguere gli attributi dai parametri con lo stesso nome (\cc|this->x = x|), per
restituire l'oggetto stesso (\cc|return *this|, come in \cc|operator=|) e per
verificare l'autoassegnamento (\cc|this == &altro|). I metodi statici non hanno
\cc|this|.
\end{orale}

\begin{orale}[D6. Come funzionano le eccezioni in C++?]
Quando una funzione incontra una situazione che non può gestire, la segnala con
\cc|throw|, lanciando un oggetto che descrive l'errore. L'esecuzione si interrompe
e risale le chiamate finché non trova un blocco \cc|try| con un \cc|catch| adatto
al tipo lanciato, che gestisce l'errore. La libreria \cc|<stdexcept>| offre classi
pronte come \cc|invalid_argument| e \cc|out_of_range|, tutte derivate da
\cc|exception|, il cui metodo \cc|what()| restituisce il messaggio. Si catturano
per reference costante. Un costruttore che lancia un'eccezione impedisce la
creazione dell'oggetto.
\end{orale}

% ---------------------------------------------------------------------
\section{Esercizi al PC}
% ---------------------------------------------------------------------

\begin{esercizio}{Rettangolo flessibile}{facile}
Classe \cc|Rettangolo| con un \textbf{unico} costruttore a parametri di default
(\cc|Rettangolo(double b = 1, double h = 1)|) che usa la lista di inizializzazione e
corregge i valori non positivi. Metodi \cc|const|: \cc|area|, \cc|perimetro|,
\cc|diagonale|, \cc|isQuadrato|. Crea un \cc|const Rettangolo| e verifica quali
metodi puoi chiamare. (È l'esercizio delle slide del corso.)
\end{esercizio}

\begin{esercizio}{Chi nasce e chi muore}{facile}
Scrivi una classe \cc|Traccia| il cui costruttore stampa \texttt{nasce <nome>} e il
cui distruttore stampa \texttt{muore <nome>}. Creane una globale, alcune locali nel
\cc|main| e in un blocco, una in una funzione e una con \cc|new|. \textbf{Prima} di
eseguire, scrivi su un foglio l'ordine che ti aspetti.
\end{esercizio}

\begin{esercizio}{Contatore di istanze}{facile}
Classe \cc|Punto| con due attributi \cc|static|: \cc|vivi| (oggetti esistenti) e
\cc|creati| (oggetti mai creati). Aggiornali in costruttore, costruttore di copia e
distruttore. Verifica cosa succede passando un \cc|Punto| per valore a una funzione.
\end{esercizio}

\begin{esercizio}{Concatenare con this}{facile}
Classe \cc|Punto| il cui costruttore ha parametri con lo stesso nome degli attributi
(usa \cc|this->|). Aggiungi \cc|setX|, \cc|setY| e \cc|sposta| che restituiscono
\cc|Punto&| per scrivere \cc|p.setX(3).setY(4).sposta(1, 1);|, e un metodo
\cc|stessoOggetto(const Punto&)| che confronta gli indirizzi.
\end{esercizio}

\begin{esercizio}{Punti e rettangoli}{medio}
Classe \cc|Punto2D| e classe \cc|Rettangolo| con due attributi \cc|Punto2D|
(\cc|topLeft| e \cc|bottomRight|) inizializzati nella lista. Metodi
\cc|contiene(Punto2D)| e \cc|contiene(const Rettangolo&)| (overloading). Deve
compilare \cc|Rettangolo B({10, 20}, {50, 10});|. (Esercizio delle slide.)
\end{esercizio}

\begin{esercizio}{Orario con eccezioni}{medio}
Classe \cc|Tempo| con ore, minuti e secondi i cui setter lanciano
\cc|invalid_argument| per valori fuori intervallo, e un costruttore che usa i setter.
Nel \cc|main| cattura un errore di un setter, un errore del costruttore (con
\cc|catch (const exception&)|) e un \cc|throw 42;| con \cc|catch (...)|.
\end{esercizio}

\begin{esercizio}{Vedere la copia superficiale}{medio}
Scrivi una classe con un attributo \cc|int*| allocato nel costruttore e
\textbf{nessun} costruttore di copia: copia un oggetto, modifica la copia e stampa
l'originale. Poi scrivi una seconda classe con la Regola dei Tre e ripeti la prova.
(Non dare un distruttore alla prima classe, altrimenti il programma farebbe un
doppio \cc|delete|: libera la memoria a mano, una volta sola.)
\end{esercizio}

\begin{esercizio}{Stringa fatta in casa}{medio}
Classe \cc|MiaStringa| che memorizza il testo in un \cc|char*| dinamico terminato da
\cc|'\0'|. Costruttore da \cc|const char*|, Regola dei Tre, \cc|size()|,
\cc|concatena(const MiaStringa&)| che restituisce una nuova stringa,
\cc|maiuscolo()| e \cc|stampa()|. Non usare \cc|string| né \cc|<cstring>|.
\end{esercizio}

\begin{esercizio}{Classe Matrice}{difficile}
Classe \cc|Matrice| con \cc|int** m|: costruttore (lancia \cc|invalid_argument| per
dimensioni non positive), distruttore, costruttore di copia, \cc|operator=| che
funziona anche tra matrici di dimensioni diverse, \cc|get|/\cc|set| che lanciano
\cc|out_of_range| e \cc|somma(const Matrice&)| che restituisce una nuova matrice
(o lancia se le dimensioni sono diverse).
\end{esercizio}

\begin{esercizio}{Pila di interi}{difficile}
Classe \cc|Pila| su array dinamico con \cc|push| (raddoppia se piena), \cc|pop| e
\cc|top| (lanciano \cc|underflow_error| se vuota), \cc|vuota|, \cc|size| e Regola
dei Tre. Verifica che svuotare una pila non svuoti la sua copia. È un'anteprima del
Modulo 12.
\end{esercizio}

\subsection*{Soluzioni}
Le soluzioni commentate, con il codice e l'output reale, sono nel PDF separato
\textbf{Soluzioni} (\texttt{soluzioni-light.pdf} / \texttt{soluzioni-dark.pdf}),
allo stesso numero di esercizio. Il codice è anche in \texttt{codice/m04/}.
